\section{Discussion and Limitations}

\paragraphb{Driver modifications.}
\brand requires modifying IO device drivers to use the new callback-based \CodeIn{unmap()} interface,
    which defers free-page operations to the IOMMU subsystem
    instead of freeing pages immediately after unmap.
In our implementation, adapting the mlx5 driver required ${\sim}350$ lines of changes,
    comparable to the driver modifications in prior work such as F\&S~\cite{Rubin2024FastSafe}.
% The changes are primarily mechanical:
%     replacing direct page-free calls with callback registration at unmap time.
However, each new driver wishing to benefit from \brand must be similarly adapted.
\optOne{} can be used without driver modifications; all existing driver will enjoy the benefits.
% We expect that a generic helper in the DMA mapping layer could reduce per-driver effort in the future.

\paragraphb{Reduced benefits at low core counts.}
\brand's one-for-many invalidations benefit from concurrent producers generating
    multiple IR-PREPs during each IR-SUB operation.
With fewer cores, IR group sizes shrink and the batching advantage diminishes.
Moreover, the pinned consumer strategy dedicates one core to IR-SUB,
    which represents a proportionally larger resource cost at low core counts.
For example, our largest gap to vIOMMU-off (8.4\%) occurs at 12 cores in the Tx-server setup.
At this core count, deferred free page dispatches releases pages in bursts
    on all producer cores at once,
    causing producer cores to contend on the buddy allocator's zone lock.
This contention slows down page reclamation and can lead to socket-buffer stalling (\S\ref{ssec:dfp}).
At higher core counts (20 and 24), there are enough producers to absorb this overhead.
We believe this gap can be further narrowed by leveraging the buddy allocator's
    existing bulk-free interfaces to amortize lock acquisition across multiple pages.

\paragraphb{Application to shadow page table-based vIOMMU}
\uvmname{} focuses on nested translation IOMMU;
however, its techniques can benefit shadow page table-based vIOMMU as well.
More specifically, shadow page table-based vIOMMU requires VM exits at IOMMU cache invalidation time at both map and unmap time
to notify the hypervisor of page table changes.

Our \optOne can be applied to both map and unmap time invalidations.
Yet, though we can not batch multiple invalidation requests into a single IR-SUB operation,
we can still do one IR-SUB but with multiple IR descriptors. This batching will reduce VM exits overhead.

Deferred free page can only be applied to unmap time invalidations.
At map time, we cannot proceed without invalidation completion because we need to 
ensure shadow page table has been correctly updated;
otherwise, handing a IOVA without shadow page tabl entry built to NIC will lead to DMA failure.


\paragraphb{Relationship to lazy mode}
Linux's lazy (deferred) mode also achieves performance benefits by batching invaladtion requests within a deferred time window.
Linux's current lazy mode is vulnerable to both confidentiality and integrity attacks~\cite{AlexCodeInjection, RajapakshaIOMMU}.

The technique of deferred free page can be applied to lazy mode to avoid the attacks.
However, with a non trivial deferred time window (10 ms by default in Linux),
    deferring free page operation may lead to much higher memory overhead in terms of active pages.
The larger the deferred time window, the more more pages will need to be freed in burst--
it will lead to huge CPU bottlenecks in the IOMMU subsystem, 
and makes it prone to the socket buffer pool exhaustion problem discussed in \S\ref{ssec:dfp}.

Note that our free page operations is only deferred for at most the duration of two rounds of IR-SUB operations (10 - 30 $\mu$s) in practice,
only less than 1\% of the deferred time window.
% From a security perspective, applying deferred free page to lazy mode still leave 
    % creating a window in which stale IOMMU cache entries permit unauthorized DMA access.
\schaic{I didn't make the claim on lazy+DFP has larger attack window for TOCTTOU attack.}

% In contrast, \brand's deferred free page defers only the \emph{page reclamation},
%     not the invalidation:
%     the invalidation still occurs before the page is returned to the allocator,
%     preserving the strict safety invariant.